\documentclass[12pt, oneside, a4paper, brazil]{abntex2}
% \usepackage[utf8]{inputenc}
\usepackage[brazil]{babel}
\usepackage{lmodern}
% \usepackage[T1]{fontenc}
\usepackage{fontspec}
\usepackage{setspace}
\usepackage{graphicx}
\usepackage{scalefnt}
\usepackage{float}
\usepackage[a4paper, left=3cm, right=2.5cm, top=3cm, bottom = 2.5cm]{geometry}
\usepackage{hyperref}
\usepackage{listings}
\usepackage{color}
\usepackage{titlesec}
%\usepackage[alf]{abntex2cite}
% \usepackage{times}
\usepackage[scaled]{helvet}
\usepackage{fancyhdr}


%%% Definição de variáveis. %%%
% ---
% Informações de dados para CAPA e FOLHA DE ROSTO
% ---
\titulo{Análise de desempenho de estruturas de dados para um indexador de arquivos}
\autor{Atílio Antônio Dadalto \\ Tiago da Cruz Santos}
\local{Vitória}
\data{2019}
\instituicao{%
  Universidade Federal do Espírito Santo
  \par Departamento de Informática}
\tipotrabalho{Relatório}
\preambulo{Relatório apresentado como requisito parcial para aprovação na disciplina de Programação I, pela Universidade Federal do Espírito Santo.}

\newcommand{\versao}{2.0}
\newcommand{\subtitulo}{Anteprojeto}


% Macros específicas do trabalho.
% (*) Inclua aqui termos que são utilizados muitas vezes e que demandam formatação especial.
% Os exemplos abaixo incluem i* (substituindo o asterisco por uma estrela) e Java com TM em superscript.
% Use sempre \xspace para que o LaTeX inclua espaço em branco após a macro somente quando necessário.
\newcommand{\istar}{\textit{i}$^\star$\xspace}
\newcommand{\java}{Java\texttrademark\xspace}
\newcommand{\latex}{\LaTeX\xspace}


%%% Configurações finais de aparência. %%%

% Altera o aspecto da cor azul.
\definecolor{blue}{RGB}{41,5,195}

% Informações do PDF.
\makeatletter
\hypersetup{
	pdftitle={\@title},
	pdfauthor={\@author},
	pdfsubject={\imprimirpreambulo},
	pdfcreator={LaTeX with abnTeX2},
	pdfkeywords={abnt}{latex}{abntex}{abntex2}{trabalho acadêmico},
	colorlinks=true,				% Colore os links (ao invés de usar caixas).
	linkcolor=blue,					% Cor dos links.
	citecolor=blue,					% Cor dos links na bibliografia.
	filecolor=magenta,				% Cor dos links de arquivo.
	urlcolor=blue,					% Cor das URLs.
	bookmarksdepth=4
}
\makeatother

% Espaçamentos entre linhas e parágrafos.
\setlength{\parindent}{1.3cm}
\setlength{\parskip}{0.2cm}



%%% Páginas iniciais do documento: capa, folha de rosto, ficha, resumo, tabelas, etc. %%%

% Compila o índice. <--- Desnecessário em Plano de Estudo.
%\makeindex

% Inicia o documento.
\begin{document}

% Retira espaço extra obsoleto entre as frases.
\frenchspacing


% Brasão da instituição.
\begin{figure}[h]
  \centering
  \includegraphics[scale=0.055]{figuras/brasao}
\end{figure}

% Capa do trabalho.
\imprimircapa
\imprimirfolhaderosto

% Lista de silgas.
% (*) Indicar as siglas utilizadas no trabalho como no exemplo abaixo.
%\begin{siglas}
    %\item [UML] Linguagem de Modelagem Unificada, do inglês \textit{Unified Modeling Language}
%\end{siglas}

% Índice de capítulos.
% \tableofcontents*


%%% Início da parte de conteúdo do documento. %%%

% Marca o início dos elementos textuais.
\clearpage
\textual

% Inclusão dos capítulos.
% (*) Para facilitar a organização, os capítulos foram divididos em arquivo separados e colocados dentro da.
% pasta capitulos/. Caso o aluno prefira trabalhar com um só arquivo, basta substituir os comandos \input
% pelos conteúdos dos arquivos que estão sendo incluídos, excluindo a pasta capitulos/ em seguida.

\input{corpo}


% Finaliza a parte no bookmark do PDF para que se inicie o bookmark na raiz e adiciona espaço de parte no sumário.
\phantompart

% Marca o fim dos elementos textuais.
\postextual

% Referências bibliográficas
\clearpage
\bibliography{bibliografia}

% Fim do documento.
\end{document}
